\documentclass[11pt,a4paper]{article}
\usepackage[utf8]{inputenc}
\usepackage[T1]{fontenc}
\usepackage[brazil]{babel}
\usepackage[margin=2.5cm]{geometry}
\usepackage{mathptmx}
\usepackage{amsmath}
\usepackage{amssymb}
\usepackage{enumitem}
\usepackage{booktabs}
\usepackage{tabularx}
\usepackage{array}
\usepackage{hyperref}
\hypersetup{colorlinks=false, pdfborder={0 0 0}}
\usepackage{fancyhdr}

\pagestyle{fancy}
\fancyhf{}
\fancyhead[L]{\small TrackFleet --- Finanças}
\fancyhead[R]{\small\thepage}
\renewcommand{\headrulewidth}{0.4pt}

\newcolumntype{L}{>{\raggedright\arraybackslash}X}
\newcolumntype{C}{>{\centering\arraybackslash}X}
\newcolumntype{R}{>{\raggedleft\arraybackslash}X}
\setlist{topsep=3pt}

% Cabeçalho padrão do projeto. No documento consolidado (trabalho_pratico_01/)
% estes três comandos são redefinidos e este mesmo arquivo vira parte do capítulo do domínio.
\newcommand{\cabecalho}[3]{%
  \begin{center}
  {\Large\bfseries DOCUMENTO DE #1 DO PROJETO}\\[2pt]
  {\large\bfseries TrackFleet --- Gestão de Rastreador de Automóveis}\\[4pt]
  Disciplina de Gerenciamento de Projetos $\cdot$ Framework PMBOK\textregistered\ 8\textsuperscript{a} Edição $\cdot$ Domínio: #2\\[2pt]
  Integrantes: Enzo Allebrand e Kerlon Ribeiro (dois GPs) $\cdot$ 4\textsuperscript{o} semestre $\cdot$ Data: #3
  \end{center}
  \vspace{2mm}\hrule\vspace{4mm}}
\newcommand{\parte}[1]{\noindent\textbf{\large #1}\par\vspace{1mm}}
\newcommand{\separador}{\par\vspace{2mm}\hrule\vspace{4mm}}

\begin{document}

\cabecalho{FINANÇAS}{Finanças}{03/09}

\noindent No PMBOK\textregistered\ 8, o domínio de desempenho Finanças cobre quatro processos: Planejar o Gerenciamento Financeiro (o manual de regras do dinheiro do projeto), Estimar Custos (quanto vai custar cada pacote de trabalho), Desenvolver o Orçamento (agregar as estimativas e as reservas em uma linha de base autorizada) e Monitorar e Controlar as Finanças (comparar o realizado com a linha de base e prever onde o projeto vai terminar). A virada de mentalidade da aula guia todo este documento: ficar dentro do orçamento é um resultado de controle; entregar o valor prometido é um resultado estratégico. Por isso o documento não para na planilha de custos --- ele fecha justificando, com ROI, VPL, TIR e payback, por que o dinheiro gasto no TrackFleet produz o resultado prometido.

\noindent Os números desta versão foram recalculados a partir do cronograma nivelado (51 dias úteis) e da reserva de contingência dimensionada pelo registro de riscos. A planilha \texttt{trackfleet\_financas.xlsx} reproduz todos os cálculos com fórmulas.

\separador

%======================================================================
\section{Planejar o Gerenciamento Financeiro: o manual de regras}\label{fin:planejar}
%======================================================================

\noindent Antes de estimar qualquer valor, este processo define como receitas e despesas do TrackFleet serão estimadas, orçadas, monitoradas e controladas. Sem ele, o resto é improviso. As entradas deste plano são o Termo de Abertura, o Documento de Escopo, o Documento de Cronograma e o Documento de Governança; a saída é o conjunto de regras abaixo.

\subsection{Bases de referência}
\begin{itemize}[itemsep=3pt]
  \item \textbf{Moeda:} real brasileiro (R\$); todos os valores deste documento estão em R\$ correntes.
  \item \textbf{Custo-hora de referência:} R\$ 50,00 por hora de trabalho de GP, usado como custo de oportunidade das horas dos dois integrantes --- o projeto não paga salário, mas contabiliza o tempo pelo valor de mercado de uma hora de desenvolvimento júnior, para que o custo real do esforço fique visível e não pareça ``de graça''.
  \item \textbf{Dia útil de trabalho:} 6 horas produtivas; logo, 1 dia-pessoa do cronograma equivale a R\$ 300,00. Essa é a ponte entre o Documento de Cronograma (durações em dias) e este documento (custos em reais).
  \item \textbf{Cadência de medição:} as finanças são medidas no mesmo check-in semanal já definido na Governança e no Cronograma --- não há reunião financeira à parte.
\end{itemize}

\subsection{Fontes de financiamento}
Cada fonte traz regras e nível de rigor diferentes. O TrackFleet combina duas das fontes vistas na aula e descarta as demais:
\begin{itemize}[itemsep=2pt]
  \item \textbf{Orçamento interno (fonte principal):} as horas de trabalho dos dois GPs, no contexto acadêmico da disciplina. É recurso próprio da dupla, sem desembolso de caixa.
  \item \textbf{Contrato/patrocínio de cliente (fonte secundária):} a transportadora parceira (patrocinador) custeia a infraestrutura do piloto --- assinatura da API do provedor de rastreamento, hospedagem e domínio.
  \item \textbf{Subvenções/grants e crowdfunding:} não se aplicam a este projeto.
\end{itemize}
\noindent Como ensina a aula, a fonte define o nível de rigor do controle: por ser financiamento interno e um patrocínio leve de parceiro, a governança financeira é leve (coerente com o modelo híbrido da Governança). A única exceção é a LGPD, que impõe controle formal de tratamento de dados independentemente da fonte.

\subsection{Classificação CapEx {\texttimes} OpEx}
Classificar cada custo na categoria certa é decidido aqui, no planejamento --- classificar errado, como alerta a aula, trava o projeto porque os caminhos de aprovação são diferentes.

\begin{center}
\begin{tabularx}{\textwidth}{l L L}
\toprule
\textbf{Categoria} & \textbf{O que é no TrackFleet} & \textbf{Itens} \\
\midrule
CapEx --- constrói o ativo & Trabalho que cria o ativo de software (a plataforma), gasto de uma vez para gerar um bem duradouro & Horas de trabalho dos GPs em todos os pacotes da EAP, incluindo arquitetura e configuração inicial da infraestrutura \\
OpEx --- roda o dia a dia & Custos contínuos de manter a plataforma no ar durante o projeto e o piloto & Assinatura mensal da API do provedor, hospedagem no VPS, domínio anual, ferramentas; certificado TLS e envio de e-mail em planos gratuitos \\
\bottomrule
\end{tabularx}
\end{center}

\noindent Um ponto concreto liga esta classificação à restrição do Termo de Abertura: o TrackFleet tem orçamento zero para hardware próprio. Na prática, o projeto elimina de propósito o maior CapEx de um sistema de rastreamento --- comprar e instalar rastreadores --- integrando-se por API a dispositivos de terceiros. Esse custo vira OpEx (assinatura da API), muito menor e sem imobilização de capital.

\subsection{Governança financeira e tolerâncias}\label{fin:tolerancias}
Quem aprova o quê segue a matriz RACI da Governança, agora aplicada ao dinheiro:
\begin{itemize}[itemsep=3pt]
  \item \textbf{Gasto do dia a dia (dentro da linha de base):} decidido pelos dois GPs por consenso, sem subir ao patrocinador.
  \item \textbf{Reserva de contingência:} está no controle dos GPs (é para riscos já conhecidos); cada uso é ligado ao risco do registro que o justificou e anotado no log de decisões.
  \item \textbf{Reserva de gerenciamento:} só é liberada pelo patrocinador (Diretoria da transportadora parceira) --- coerente com a aula, que deixa essa reserva com a alta gerência, e com o escalonamento já definido na Governança.
\end{itemize}

\noindent As tolerâncias abaixo são os gatilhos que transformam um número em uma ação --- porque, como diz a aula, relatório colorido parado não é controle:
\begin{center}
\begin{tabularx}{\textwidth}{L L}
\toprule
\textbf{Gatilho} & \textbf{Ação} \\
\midrule
CPI ou SPI abaixo de 0,90 em um check-in & Ação corretiva imediata dos GPs, focada primeiro no caminho crítico \\
Projeção de término (EAC) acima da linha de base, ameaçando a reserva de gerenciamento & Revisão formal e alerta ao patrocinador; acima da reserva, escalonamento e re-baseline \\
Consumo da reserva de contingência acima do avanço físico do projeto & Reavaliação das estimativas restantes e do registro de riscos no check-in seguinte \\
\bottomrule
\end{tabularx}
\end{center}

\separador

%======================================================================
\section{Estimar Custos: quanto vai custar cada pacote}\label{fin:estimar}
%======================================================================

\noindent Estimar custos é uma aproximação periódica do custo dos recursos necessários ao trabalho. A base é a EAP do Documento de Escopo e as atividades do Documento de Cronograma: o custo de cada pacote é a soma dos dias-pessoa das suas atividades, multiplicada por R\$ 300,00. Como na aula, a técnica usada em cada pacote segue a natureza do trabalho.

\subsection{Técnica por pacote}
\begin{itemize}[itemsep=3pt]
  \item \textbf{Análoga:} a gestão do projeto (1.1), que não é uma atividade do cronograma e sim um esforço contínuo (check-ins, relatórios e documentos de cada domínio), é estimada por comparação com o tempo que a dupla já gastou nos documentos das aulas; requisitos (1.2) e encerramento (1.5) seguem as estimativas análogas do Cronograma.
  \item \textbf{Paramétrica:} o Frontend (1.3.3) usa as estimativas paramétricas do Cronograma --- número de telas e de relatórios do catálogo do Documento de Escopo multiplicado pelo tempo médio por unidade.
  \item \textbf{Especialista e bottom-up:} pacotes pequenos e bem conhecidos (autenticação e LGPD, infraestrutura) são estimados somando suas atividades diretamente.
  \item \textbf{Três pontos (PERT):} os três pacotes que contêm as atividades estimadas em três pontos no Cronograma, detalhados na Seção~\ref{fin:pert}.
\end{itemize}

\subsection{Estimativa bottom-up por pacote}\label{fin:pacotes}
Cada pacote é convertido em dias-pessoa a partir das atividades do Cronograma e, pelo custo de R\$ 300,00 por dia-pessoa, em reais. No piloto (A18), cada GP dedica um terço da jornada; A19 e A20 usam os outros dois terços dos seus dias.

\begin{center}
\small
\begin{tabularx}{\textwidth}{l L c c r}
\toprule
\textbf{Pacote EAP} & \textbf{Atividades do Cronograma} & \textbf{Técnica} & \textbf{Dias-pessoa} & \textbf{Custo (R\$)} \\
\midrule
1.1 Gestão do Projeto & Esforço contínuo de gestão (fora da rede) & Análoga & 12,00 & 3.600 \\
1.2 Requisitos e Arquitetura & A1 (dois GPs), A2, A3, A4 & Análoga & 13,00 & 3.900 \\
1.3.1 Backend, API e ingestão & A5, A6, A7 & Três pontos & 15,00 & 4.500 \\
1.3.2 Motor de alertas & A8 & Três pontos & 9,00 & 2.700 \\
1.3.3 Frontend & A11, A12, A13, A14 & Paramétrica & 17,00 & 5.100 \\
1.3.4 Autenticação e LGPD & A9, A10 & Bottom-up & 5,00 & 1.500 \\
1.3.5 Infraestrutura e implantação & A21 & Especialista & 2,00 & 600 \\
1.4 Testes e Homologação & A15, A16, A17, A18 (1/3 de cada GP) & Três pontos & 17,33 & 5.200 \\
1.5 Treinamento e Encerramento & A19 (dois GPs, 2/3), A20 (2/3) & Análoga & 3,33 & 1.000 \\
\midrule
\textbf{Subtotal mão de obra (CapEx)} & & & \textbf{93,67} & \textbf{28.100} \\
Infraestrutura do período (OpEx) & API, VPS, domínio, ferramentas & Bottom-up & --- & 700 \\
\midrule
\textbf{Total das estimativas} & & & & \textbf{28.800} \\
\bottomrule
\end{tabularx}
\end{center}

\subsection{Estimativa de três pontos dos pacotes de risco}\label{fin:pert}
Os três pacotes mais incertos usam os cenários otimista (O), mais provável (M) e pessimista (P) das atividades do Cronograma, convertidos em reais pela mesma fórmula da distribuição Beta:
\[
C_E = \frac{O + 4M + P}{6}
\]

\begin{center}
\begin{tabularx}{\textwidth}{L r r r r r}
\toprule
\textbf{Pacote} & \textbf{O} & \textbf{M} & \textbf{P} & $C_E$ & \textbf{Na linha de base} \\
\midrule
1.3.1 Backend, API e ingestão & 2.700 & 3.900 & 6.600 & 4.150 & 4.500 \\
1.3.2 Motor de alertas & 1.500 & 2.400 & 4.500 & 2.600 & 2.700 \\
1.4 Testes e Homologação & 4.400 & 5.000 & 6.000 & 5.067 & 5.200 \\
\bottomrule
\end{tabularx}
\end{center}

\noindent A coluna ``na linha de base'' é maior que $C_E$ porque o Cronograma arredonda cada $t_E$ para cima (dia inteiro de trabalho). A leitura desses números: o Backend (1.3.1) e o Motor de alertas (1.3.2) carregam a maior distância entre otimista e pessimista --- o Backend depende de uma API externa fora do controle da dupla, e o Motor de alertas reúne quatro lógicas de disparo diferentes na mesma frente. Essa incerteza é tratada no registro de riscos, que dimensiona a reserva de contingência.

\separador

%======================================================================
\section{Desenvolver o Orçamento: da estimativa à linha de base}\label{fin:orcamento}
%======================================================================

\noindent Aqui as estimativas individuais são agregadas e as reservas são adicionadas, formando a linha de base de custos autorizada. Este é o ponto onde mais se erra na prova, como diz a aula: orçamento não é a mesma coisa que linha de base, e as duas reservas entram em lugares diferentes.

\subsection{As duas reservas}
\begin{center}
\begin{tabularx}{\textwidth}{L L L}
\toprule
\textbf{Reserva} & \textbf{Cobre} & \textbf{Onde fica} \\
\midrule
Contingência --- riscos conhecidos & Os riscos do registro de riscos, pelo valor monetário esperado (VME) líquido das ameaças e oportunidades & DENTRO da linha de base; já entra no orçamento inicial e é controlada pelos GPs \\
Gerenciamento --- riscos desconhecidos & Trabalho imprevisto que ninguém mapeou & FORA da linha de base; entra no orçamento total por cima dela e só é liberada pelo patrocinador \\
\bottomrule
\end{tabularx}
\end{center}

\noindent O mnemônico da aula vale para o TrackFleet: contingência para o conhecido.

\subsection{Composição do orçamento}
\begin{center}
\begin{tabularx}{\textwidth}{L r}
\toprule
\textbf{Componente} & \textbf{Valor (R\$)} \\
\midrule
Estimativas dos pacotes (mão de obra) & 28.100 \\
Infraestrutura do período (OpEx) & 700 \\
\textbf{Subtotal das estimativas} & \textbf{28.800} \\
(+) Reserva de contingência (VME líquido do registro de riscos) & 3.970 \\
\textbf{Linha de base de custos (BAC)} & \textbf{32.770} \\
(+) Reserva de gerenciamento (5\% da linha de base) & 1.639 \\
\textbf{Orçamento total do projeto} & \textbf{34.409} \\
\bottomrule
\end{tabularx}
\end{center}

\noindent A linha de base de custos (BAC) é R\$ 32.770: é contra ela que o realizado será medido na Seção~\ref{fin:controlar}. O orçamento total (R\$ 34.409) é maior porque inclui a reserva de gerenciamento --- que existe, mas não faz parte da base de comparação. A reserva de contingência (R\$ 3.970, cerca de 14\% das estimativas) não é um percentual chutado: é a soma dos VMEs do registro de riscos. A linha de base é aprovada pelo patrocinador, na mesma lógica das linhas de base de escopo e de cronograma.

\subsection{Orçamento no tempo: a curva S}\label{fin:curvas}
A linha de base é distribuída no tempo sobre as 11 semanas do Cronograma: cada semana recebe o custo dos dias-pessoa trabalhados nela, uma fração igual da gestão do projeto e do OpEx, e a reserva de contingência proporcional ao custo.

\begin{center}
\footnotesize
\setlength{\tabcolsep}{3pt}
\begin{tabularx}{\textwidth}{l *{11}{C}}
\toprule
\textbf{Semana} & \textbf{S1} & \textbf{S2} & \textbf{S3} & \textbf{S4} & \textbf{S5} & \textbf{S6} & \textbf{S7} & \textbf{S8} & \textbf{S9} & \textbf{S10} & \textbf{S11} \\
\midrule
\% acumulado & 11,8 & 20,4 & 32,2 & 44,0 & 55,7 & 67,5 & 77,2 & 84,8 & 89,6 & 93,8 & 100 \\
Acumulado (R\$) & 3.858 & 6.693 & 10.551 & 14.409 & 18.268 & 22.126 & 25.302 & 27.795 & 29.377 & 30.732 & 32.770 \\
\bottomrule
\end{tabularx}
\end{center}

\noindent A curva tem o miolo íngreme e o fim achatado previstos na aula --- no piloto (S9 a S11) a equipe trabalha só um terço da jornada. O arranque, porém, não é lento: com uma equipe fixa de duas pessoas alocadas desde o primeiro dia, o gasto já começa no ritmo de cruzeiro. Essa curva é o Valor Planejado (PV) ao longo do tempo e é a régua do exemplo resolvido da próxima seção.

\separador

%======================================================================
\section{Monitorar e Controlar as Finanças: o valor agregado}\label{fin:controlar}
%======================================================================

\noindent Este processo supervisiona a saúde financeira: acompanha gastos, ajusta previsões e aplica correções. A ferramenta central é o Gerenciamento do Valor Agregado (EVM), começado desde o dia 1 --- porque, como avisa a aula, fazer EVM só aos 60\% do projeto é inútil: a janela de intervenção já fechou.

\subsection{Os quatro valores-base}
\begin{center}
\begin{tabularx}{\textwidth}{l L L}
\toprule
\textbf{Sigla} & \textbf{Nome} & \textbf{O que é no TrackFleet} \\
\midrule
BAC & Orçamento no término & A linha de base de custos: R\$ 32.770 \\
PV & Valor Planejado & Custo orçado do trabalho programado até a data (a curva S da Seção~\ref{fin:curvas}) \\
EV & Valor Agregado & Custo orçado do trabalho realmente concluído: BAC $\times$ \% concluído \\
AC & Custo Real & O que de fato foi consumido para entregar esse trabalho \\
\bottomrule
\end{tabularx}
\end{center}

\subsection{Variações, índices e previsões}\label{fin:formulas}
\begin{center}
\begin{tabularx}{\textwidth}{L L L}
\toprule
\textbf{Indicador} & \textbf{Fórmula} & \textbf{Leitura} \\
\midrule
CV --- Variação de Custo & $EV - AC$ & Positivo: abaixo do orçamento. Negativo: acima \\
SV --- Variação de Prazo & $EV - PV$ & Positivo: adiantado. Negativo: atrasado \\
CPI --- Índice de Custo & $EV \div AC$ & Acima de 1,0 é bom; abaixo é ruim \\
SPI --- Índice de Prazo & $EV \div PV$ & Acima de 1,0 é bom; abaixo é ruim \\
EAC --- Estimativa no Término & $BAC \div CPI$ & Custo final previsto se a tendência atual continuar \\
ETC --- Estimativa para Terminar & $EAC - AC$ & Quanto ainda falta gastar \\
VAC --- Variação no Término & $BAC - EAC$ & Negativo: projeto termina acima da linha de base \\
TCPI --- Índice para Terminar & $(BAC - EV) \div (BAC - AC)$ & Eficiência exigida do trabalho restante; acima de 1,0 é difícil \\
\bottomrule
\end{tabularx}
\end{center}

\noindent Duas armadilhas da aula ficam registradas aqui: termina em ``Variação'' subtrai, termina em ``Índice'' divide; e SV/SPI medem prazo, mas em unidades de dinheiro (o valor em R\$ do trabalho), nunca em dias.

\subsection{Exemplo resolvido: check-in do fim da Semana 6}\label{fin:exemplo}
Cenário de acompanhamento no corte da Semana 6 (09/10), quando a curva S prevê 67,5\% da linha de base. Suponha que, na medição, o trabalho realmente concluído seja 62\% e o custo real acumulado seja R\$ 21.200:

\begin{center}
\begin{tabularx}{\textwidth}{L L L}
\toprule
\textbf{Valor} & \textbf{Cálculo} & \textbf{Resultado} \\
\midrule
PV & curva S no fim da S6 & R\$ 22.126 \\
EV & $32.770 \times 62\%$ & R\$ 20.317 \\
AC & medido no check-in & R\$ 21.200 \\
CV & $20.317 - 21.200$ & $-$R\$ 883 (acima do orçamento) \\
SV & $20.317 - 22.126$ & $-$R\$ 1.809 (atrasado) \\
CPI & $20.317 \div 21.200$ & 0,96 \\
SPI & $20.317 \div 22.126$ & 0,92 \\
EAC & $32.770 \div (20.317 \div 21.200)$ & R\$ 34.194 \\
ETC & $34.194 - 21.200$ & R\$ 12.994 \\
VAC & $32.770 - 34.194$ & $-$R\$ 1.424 \\
TCPI & $(32.770 - 20.317) \div (32.770 - 21.200)$ & 1,08 \\
\bottomrule
\end{tabularx}
\end{center}

\noindent \textbf{Veredito e decisão:} no ritmo atual o TrackFleet está levemente acima do orçamento (CPI 0,96) e atrasado (SPI 0,92), e a projeção de término (EAC R\$ 34.194) fura a linha de base em R\$ 1.424 --- o equivalente a 87\% da reserva de gerenciamento (R\$ 1.639). Pela tolerância da Seção~\ref{fin:tolerancias}, isso ainda não estoura o orçamento total, então não é escalonamento imediato, mas dispara revisão formal e coloca o patrocinador em alerta. A ação corretiva se concentra no maior risco do caminho crítico, o Motor de alertas (atividade A8 do Cronograma). O TCPI de 1,08 diz o preço disso: o trabalho restante precisa ser feito cerca de 8\% mais eficiente do que o planejado --- exigente, mas possível. E, seguindo a dica do professor, a decisão olha a tendência do CPI ao longo dos check-ins, não a foto de uma semana isolada.

\separador

%======================================================================
\section{O Valor: por que o investimento se justifica}\label{fin:valor}
%======================================================================

\noindent Cumprir o orçamento não é sucesso financeiro. Esta seção responde à pergunta executiva da aula --- ``o dinheiro gasto produziu o resultado prometido?'' --- medindo o retorno do TrackFleet para o patrocinador. O valor, como reconhece o padrão, não é só financeiro: é tratado nas duas dimensões.

\subsection{Valor tangível: o business case}
As premissas abaixo são estimativas de business case do ponto de vista da transportadora parceira; são deliberadamente conservadoras e assumidamente frágeis (por isso a análise de sensibilidade adiante). Elas não se confundem com o ``indicador de economia em reais'' que o Documento de Escopo deixou de fora do produto --- aqui o cálculo é do projeto, não uma tela da plataforma.

\begin{itemize}[itemsep=2pt]
  \item \textbf{Investimento inicial:} orçamento total do projeto, R\$ 34.409. É uma visão conservadora: inclui o custo de oportunidade das horas da dupla, que a transportadora não desembolsa.
  \item \textbf{Ganho anual estimado:} R\$ 24.000/ano, pela redução de manutenção corretiva, menos multas por documento vencido e menor uso indevido dos veículos.
  \item \textbf{Custo operacional anual (pós-projeto):} R\$ 3.600/ano de OpEx (API, hospedagem e domínio).
  \item \textbf{Ganho líquido anual:} R\$ 24.000 $-$ R\$ 3.600 $=$ R\$ 20.400/ano.
  \item \textbf{Taxa de desconto (custo de capital):} 12\% ao ano.
\end{itemize}

\begin{center}
\begin{tabularx}{\textwidth}{L L L}
\toprule
\textbf{Métrica} & \textbf{Cálculo} & \textbf{Resultado} \\
\midrule
Payback & $34.409 \div 20.400$ & $\approx$ 1,7 ano ($\approx$ 20 meses) \\
VPL (3 anos) & $-34.409 + \sum_{t=1}^{3} \dfrac{20.400}{1{,}12^{t}}$ & $\approx +$R\$ 14.588 \\
TIR (3 anos) & taxa que zera o VPL & $\approx$ 35\% ao ano \\
ROI (3 anos) & $(3 \times 20.400 - 34.409) \div 34.409$ & $\approx$ 78\% \\
\bottomrule
\end{tabularx}
\end{center}

\noindent Leitura: o VPL é positivo (o projeto cria valor mesmo descontando o dinheiro no tempo), a TIR de cerca de 35\% supera com folga o custo de capital de 12\% (logo, aprova-se o investimento) e o payback fica abaixo de dois anos.

\noindent Sensibilidade: mesmo que o ganho anual seja 30\% menor que o previsto (R\$ 16.800 em vez de R\$ 24.000), o ganho líquido cai para R\$ 13.200/ano e o payback estica para cerca de 2,6 anos. O VPL de 3 anos fica negativo (cerca de $-$R\$ 2.705) e volta a positivo em um horizonte de 4 anos (cerca de $+$R\$ 5.684): a decisão de investir se sustenta, com margem mais estreita e dependendo de a plataforma ser usada por mais tempo.

\subsection{Valor intangível}
Nem todo valor cabe na planilha. O TrackFleet entrega também:
\begin{itemize}[itemsep=2pt]
  \item \textbf{Conformidade regulatória:} adequação à LGPD no tratamento de dados de localização e de motoristas --- valor obrigatório, não negociável (VBS do Documento de Escopo).
  \item \textbf{Segurança:} maior proteção de veículos e motoristas com o monitoramento e os alertas.
  \item \textbf{Decisão baseada em dados:} gestão da frota apoiada em dados reais, não em achismo --- o valor esperado do Termo de Abertura.
  \item \textbf{Satisfação do cliente:} uma solução acessível para um segmento hoje mal atendido pelo custo elevado das plataformas corporativas.
  \item \textbf{Aprendizado da dupla:} domínio prático do PMBOK\textregistered\ 8 aplicado de ponta a ponta, valor acadêmico do projeto.
\end{itemize}

\subsection{Alinhamento com o valor do produto}
O business case reforça a priorização já registrada na Value Breakdown Structure do Documento de Escopo: Monitoramento em tempo real (E1) e Alertas automáticos (E2), que concentram 70\% do valor, são exatamente os entregáveis que produzem a economia de manutenção corretiva e de uso indevido usada no cálculo acima. Se o prazo do semestre apertar, cortar relatórios avançados preserva o retorno; cortar E1 ou E2 destruiria o business case.

\separador

%======================================================================
\section{Armadilhas Financeiras e como o TrackFleet as evita}\label{fin:armadilhas}
%======================================================================

\noindent As pegadinhas da aula não estão nos números, e sim no comportamento diante deles. Cada uma tem uma decisão concreta no projeto:

\begin{itemize}[itemsep=5pt]
  \item \textbf{Confundir linha de base com orçamento:} a reserva de gerenciamento (R\$ 1.639) fica de fora da linha de base (R\$ 32.770) e só entra no orçamento total (R\$ 34.409). A comparação do EVM é sempre contra a linha de base, nunca contra o orçamento total.
  \item \textbf{Trocar índice por variação:} a tabela de fórmulas (Seção~\ref{fin:formulas}) registra que termina em ``Índice'' divide e termina em ``Variação'' subtrai --- e que SV/SPI medem prazo em reais, não em dias.
  \item \textbf{Inverter a leitura do TCPI:} para CPI e SPI, abaixo de 1,0 é ruim; para o TCPI, o oposto --- acima de 1,0 significa que o trabalho restante precisa ser mais eficiente, como no exemplo da Seção~\ref{fin:exemplo} (1,08).
  \item \textbf{Achar que orçamento cumprido é sucesso:} por isso este documento não para no controle de custos --- a Seção~\ref{fin:valor} mede o valor entregue (ROI, VPL, TIR, payback), que é o resultado estratégico.
  \item \textbf{Começar o EVM tarde demais:} a medição começa no dia 1, no mesmo check-in semanal, para que o desvio apareça enquanto a janela de intervenção ainda está aberta.
  \item \textbf{Relatório sem narrativa:} cada indicador do exemplo resolvido vem com a decisão que ele dispara (revisão, ação corretiva, alerta ao patrocinador) --- dado sem decisão não é controle.
  \item \textbf{Escolher o EAC errado:} como o desvio do TrackFleet é uma variação típica (tendência de custo), o EAC usa $BAC \div CPI$. A fórmula de evento único, $AC + (BAC - EV)$, só valeria para um desvio isolado que não se repete; e a fórmula pessimista, $AC + (BAC - EV) \div (CPI \times SPI)$, fica para quando custo e prazo pesam juntos sobre o restante do projeto.
\end{itemize}

\end{document}
